% Einheit 2 von 4 – Prozentwert berechnen (Katalog Einheit 3), Kl. 7, schwach
\clearpage
\hypertarget{e2}{}
\einheitenkopf*{Einheit 2 von 4 · Prozentwert berechnen}
\zweigzeile{Hier lernst du auszurechnen, wie viel ein Prozentanteil von einer Größe ist – in Euro, kg oder Stück · neu in diesem Jahr · P10 · baut auf: Prozente als Anteile, Bruchteil einer Größe, Kommazahlen malnehmen, Dreisatz (Kennst du schon)}
\setcounter{aufgabe}{21}

\begin{aufgabe}{Ich kann am Streifen ausrechnen, wie viel 50\,\%, 25\,\% oder 10\,\% von einer Größe sind. Der ganze Streifen ist das Ganze. Rechne aus, wie viel der gefärbte Teil ist.}
\swb{\small Der Streifen sind 30\,€.\\[2pt]\beispiel{50\,\% &= \text{die Hälfte} \\ 30\,€ : 2 &= 15\,€}}{\streifenfrage{50}{$0\,€$}{$30\,€$}}
\anw{Der Streifen sind jetzt immer 60\,kg.}
\swz{100\,\% von 60\,kg}{\streifenfrage{100}{$0$\,kg}{$60$\,kg}}
\swz{50\,\% von 60\,kg}{\streifenfrage{50}{$0$\,kg}{$60$\,kg}}
\swz{25\,\% von 60\,kg}{\streifenfrage[20]{25}{$0$\,kg}{$60$\,kg}}
\swz{10\,\% von 60\,kg}{\streifenfrage{10}{$0$\,kg}{$60$\,kg}}
\swfrage{Was bleibt gleich, was ändert sich?}
\end{aufgabe}

\begin{aufgabe}{Ich kann über 1\,\% ausrechnen, wie viel ein Prozentanteil ist. Rechne zuerst 1\,\% aus: Ganzes : 100. Nimm dann mal die Prozentzahl. Das Ergebnis heißt Prozentwert.}
\swb{\small 3\,\% von 200\,€\\[2pt]\beispiel{1\,\% &= 200\,€ : 100 = 2\,€ \\ 3\,\% &= 2\,€ \cdot 3 = 6\,€}}{\streifenfrage[0]{3}{$0\,€$}{$200\,€$}}
\swz{7\,\% von 300\,€}{\streifenfrage[0]{7}{$0\,€$}{$300\,€$}}
\anw{Rechne erst 10\,\% aus (Ganzes : 10). Nimm dann mal 3 oder rechne die Hälfte dazu.}
\swz{30\,\% von 90\,€}{\streifenfrage{30}{$0\,€$}{$90\,€$}}
\swz{15\,\% von 60\,€}{\streifenfrage[20]{15}{$0\,€$}{$60\,€$}}
\end{aufgabe}

\begin{aufgabe}{Ich kann einen Prozentanteil mit einer Kommazahl ausrechnen. Schreibe die Prozentzahl als Kommazahl und nimm sie mal das Ganze.}
\swb{\small 40\,\% von 35\,€\\[2pt]\beispiel{40\,\% &= 0{,}4 \\ 0{,}4 \cdot 35\,€ &= 14\,€}}{\streifenfrage{40}{$0\,€$}{$35\,€$}}
\swz{60\,\% von 45\,kg}{\streifenfrage{60}{$0$\,kg}{$45$\,kg}}
\swz{30\,\% von 2,50\,€}{\streifenfrage{30}{$0\,€$}{$2{,}50\,€$}}
\end{aufgabe}

\begin{aufgabe}{Ich kann in einem Text ausrechnen, wie viel Euro das sind. Lies genau, was gefragt ist.}
\anw{Eine Jacke kostet 60\,€. Es gibt 25\,\% Rabatt.}
\swz{Wie viel Euro spart man?}{\streifenfrage[20]{25}{$0\,€$}{$60\,€$}}
\swz{Wie viel kostet die Jacke dann noch?}{\streifenfrage[20]{75}{$0\,€$}{$60\,€$}}
\anw{Im letzten Jahr hatte der Sportverein 60 Mitglieder. Jetzt sind es 150\,\% davon.}
\swz{Wie viele Mitglieder hat der Verein jetzt?}{\zahlenstrahl[xmin=0,xmax=1.55,xstep=0.1,karo=0.6]{1/100\,\%, 1.5/150\,\%}}
\end{aufgabe}

\begin{aufgabe}{Ich kann in einem Text ausrechnen, wie viel Euro das sind – weiter.}
\swz{Ein Fahrrad kostet ohne Steuer 300\,€. Dazu kommen 19\,\% Mehrwertsteuer. Wie viel Euro Steuer sind das?}{\streifenfrage[0]{19}{$0\,€$}{$300\,€$}}
\begin{teile}
\teil Ein Tablet kostet 360\,€. Im Angebot gibt es 15\,\% Rabatt. Wie viel Euro spart man? Kreuze an. (P10 2021) \\
\kreuz{15\,€}\quad\kreuz{54\,€}\quad\kreuz{306\,€}\quad\kreuz{345\,€}
\end{teile}
\end{aufgabe}

\begin{aufgabe}{Ich finde den Fehler beim Ausrechnen eines Prozentwerts. Pia soll 40\,\% von 25\,€ ausrechnen. Sie rechnet so:}
\rechnung{0{,}04 \cdot 25\,€ &= 1\,€}
\swz{Finde den Fehler, benenne ihn und korrigiere die Rechnung.}{\streifenfrage{40}{$0\,€$}{$25\,€$}}
\swz{Rechne 80\,\% von 35\,€ mit einer Kommazahl aus.}{\streifenfrage{80}{$0\,€$}{$35\,€$}}
\end{aufgabe}

\begin{aufgabe}{Ich finde den Fehler beim Ausrechnen eines Prozentwerts – weiter. Ein Roller kostet 70\,€. Es gibt 10\,\% Rabatt. Gefragt ist: Wie viel Euro spart man? Max rechnet so:}
\rechnung{10\,\% \text{ von } 70\,€ &= 7\,€ \\ 70\,€ - 7\,€ &= 63\,€ \\ \text{Antwort: Man spart } &63\,€.}
\swz{Finde den Fehler, benenne ihn und korrigiere die Antwort.}{\streifenfrage{10}{$0\,€$}{$70\,€$}}
\swz{Ein Spiel kostet 80\,€. Es gibt 15\,\% Rabatt. Wie viel Euro spart man?}{\streifenfrage[20]{15}{$0\,€$}{$80\,€$}}
\end{aufgabe}

\begin{aufgabe}{Ich kann begründen, warum ich zuerst 1\,\% ausrechne.}
\swz{Warum rechnet man bei 9\,\% von 400\,€ zuerst 400\,€ : 100? Begründe.}{\streifen[0]{9}{$0\,€$}{$400\,€$}}
\swz{Kann 5\,\% von 200\,€ mehr sein als 10\,\% von 90\,€? Entscheide und begründe.}{\streifen[20]{5}{$0\,€$}{$200\,€$}\par\medskip\streifen{10}{$0\,€$}{$90\,€$}}
\end{aufgabe}

\begin{aufgabe}{Ich kann entscheiden, welches Angebot günstiger ist. Dieselben Kopfhörer kosten in zwei Läden 80\,€. Laden Nord gibt 25\,\% Rabatt. Laden Süd macht sie 15\,€ billiger.}
\swz{Wie viel Euro spart man bei Laden Nord?}{\streifenfrage[20]{25}{$0\,€$}{$80\,€$}}
\swz{Wie viel kosten die Kopfhörer in jedem Laden?}{\streifenfrage[20]{75}{$0\,€$}{$80\,€$}}
\swz{In welchem Laden kauft man günstiger? Begründe.}{}
\end{aufgabe}

\uebersichtskasten{Prozentwert: erst 1\,\% (Ganzes : 100), dann mal p – oder p als Dezimalzahl mal Ganzes. \\ 18\,\% von 40\,kg: \ 1\,\% $= 0{,}4$\,kg $\to$ 18\,\% $= 7{,}2$\,kg \qquad $0{,}18 \cdot 40 = 7{,}2$}
